\documentclass[ ../main.tex]{subfiles}
\providecommand{\mainx}{..}
\begin{document}
\section{Implementation}
\label{sec:impl}
The \gls{gls-phf} construction (a minimal perfect hash followed by a per-slot
fingerprint, \cref{sec:phf}) is realized directly by mature minimal-perfect-hash
libraries: the perfect hash is built once over the static key set, and the
fingerprint matrix $\matrx{V}$ is a flat array of $M$-bit words indexed by the
hash. The reference implementation is the \texttt{maph} library~\cite{maph}.
Practical deployments use established minimal-perfect-hash constructions such as
PHOBIC~\cite{phobic} and hash, displace, and compress (CHD)~\cite{chd}, which
build a perfect hash over $m$ keys in linear time and a few bits per key, so the
dominant cost of the \gls{gls-phf} (the perfect hash) is supplied off the shelf.

Because the \gls{gls-phf} stores each approximate set as a perfect hash paired
with a fingerprint array, set-theoretic compositions over non-enumerable
operands are computed at query time rather than by rebuilding. Given two
approximate sets $\ASet{S}[1]$ and $\ASet{S}[2]$, the union is the disjunction of
the two membership tests,
\begin{equation}
\SetContains[x][\left(\SetUnion[\ASet{S}[1]][\ASet{S}[2]]\right)] \equiv \left(\SetContains[x][\ASet{S}[1]]\right) \lor \left(\SetContains[x][\ASet{S}[2]]\right)\,,
\end{equation}
and any set-theoretic composition reduces to unions and complements over the
stored operands. The space and induced-rate consequences of these
stored-and-composed operations are analyzed in \cref{sec:adt}.
\end{document}
